\begin{frame}{Where we stand: NORA/MIM + the W3C stack}
	\small
	MIM (\emph{Metamodel voor Informatiemodellen}) — NORA's Informatielaag —
	splits modelling into five types; in this deck they are L0--L4, and we cover L0--L3.
	\vspace{0.3em}
	\footnotesize
	\setlength{\tabcolsep}{3pt}
	\begin{center}
	\begin{tabular}{@{}p{2.7cm}p{3.2cm}p{4.4cm}@{}}
		\toprule
		\textbf{MIM type} & \textbf{W3C / formalism} & \textbf{In this deck}\\
		\midrule
		Model van begrippen & \textbf{SKOS} vocabulaire & \textbf{L0} Concept hierarchy\\
		\midrule
		Conceptueel informatiemodel & \textbf{RDFS} + \textbf{OWL} & \textbf{L1} Information Model\\
		Logisch informatiemodel & \textbf{SHACL} constraints & \textbf{L2} Profile\\
		Fysiek / technisch gegevensmodel & \textbf{schema}: Avro, JSON-Schema, SQL & \textbf{L3} Serializations\\
		\midrule
		\color{gray}Registraties & \color{gray}instance = record, row, JSON doc & \color{gray}\textbf{L4} runtime data --- \emph{mostly out of scope}\\
		\bottomrule
	\end{tabular}
	\end{center}
	\vspace{0.3em}
	\small
	The \emph{model van begrippen} is our L0 — a \emph{SKOS} vocabulary:
	``\emph{een vocabulaire is geen informatiemodel}'' —
	terms we \emph{reuse}, not design.
\end{frame}
